\part{Seeing what happens: VGA text mode}

\chapter{Video memory}

\begin{goals}
  \item How the VGA card's text mode works
  \item What each byte at \hexa{B8000} means
  \item How to compute the address of any position on screen
\end{goals}

\section{An 80 by 25 grid}

When the PC boots, the graphics card is in \term{text mode}: it does not display pixels, it
displays characters. The screen is a grid of 80 columns by 25 rows, that is 2,000 cells.

Each cell is described by \hi{two bytes} in memory:

\begin{center}
\begin{tikzpicture}[font=\sffamily\small]
  \draw[draw=grayborder,fill=accentlight] (0,0) rectangle (3.6,1);
  \node at (1.8,0.5) {ASCII code};
  \node[font=\ttfamily\tiny,text=gray] at (1.8,-0.3) {even (low) byte};

  \draw[draw=grayborder,fill=amberlight] (3.6,0) rectangle (7.2,1);
  \node at (5.4,0.5) {colour attribute};
  \node[font=\ttfamily\tiny,text=gray] at (5.4,-0.3) {odd (high) byte};

  \draw[decorate,decoration={brace,amplitude=5pt},gray] (0,1.15) -- (7.2,1.15)
    node[midway,above,font=\sffamily\scriptsize,text=gray] {one screen cell = 16 bits};
\end{tikzpicture}
\end{center}

And the attribute byte splits again into two 4-bit halves:

\begin{center}
\begin{tikzpicture}[font=\sffamily\small]
  \draw[draw=grayborder,fill=redlight] (0,0) rectangle (3.0,0.9);
  \node at (1.5,0.45) {background (4 bits)};
  \draw[draw=grayborder,fill=greenlight] (3.0,0) rectangle (6.0,0.9);
  \node at (4.5,0.45) {foreground (4 bits)};
  \node[font=\ttfamily\tiny,text=gray] at (0.4,-0.3) {bit 7};
  \node[font=\ttfamily\tiny,text=gray] at (5.6,-0.3) {bit 0};
\end{tikzpicture}
\end{center}

Four bits give 16 possible colours, the same ones PCs had in the 1980s:

\begin{center}
\begin{tabularx}{0.95\textwidth}{@{}llllX@{}}
\toprule
\textbf{Value} & \textbf{Colour} & \textbf{Value} & \textbf{Colour} & \\
\midrule
0 & Black      & 8  & Dark grey \\
1 & Blue       & 9  & Light blue \\
2 & Green      & 10 & Light green \\
3 & Cyan       & 11 & Light cyan \\
4 & Red        & 12 & Light red \\
5 & Magenta    & 13 & Light magenta \\
6 & Brown      & 14 & Yellow \\
7 & Light grey & 15 & White \\
\bottomrule
\end{tabularx}
\end{center}

\filetag{lytlnybl/kernel/drivers/vga\_text.c}
\begin{lstlisting}[style=c]
void vga_text_set_color(vga_text* terminal, vga_color f, vga_color b) {
    terminal->color = ((uint8_t)b << 4) | (uint8_t)f;
}
\end{lstlisting}

One line that sums up all of Chapter 2: shift the background four bits left to place it in the high
half, then combine with OR.

\section{Computing a cell's address}

The 2,000 cells sit in memory one after another, row by row. The cell at row $r$ and column $c$ is
at index:

\begin{center}
\cod{index = row $\times$ 80 + column}
\end{center}

And since each cell occupies 2 bytes, its memory address is \cod{0xB8000 + index $\times$ 2}. In C
this simplifies a lot by declaring the buffer as a pointer to \cod{uint16\_t}, because pointer
arithmetic then multiplies by 2 for you:

\filetag{lytlnybl/kernel/drivers/vga\_text.c}
\begin{lstlisting}[style=c]
void vga_text_putchar(vga_text* terminal, char c) {
    uint8_t color = terminal->color;
    size_t index = terminal->row * terminal->width + terminal->column;
    uint16_t entry = ((uint16_t)color << 8) | (uint16_t)c;
    terminal->buffer[index] = entry;
}
\end{lstlisting}

Four lines and you have screen output. Compare that with what it would cost on a modern graphics
system: no driver, no compositor, no fonts. You write two bytes and a letter appears.

\begin{tryit}[Write to the screen by hand]
This is exactly what the kernel does the moment it enters C, before any driver exists:
\begin{lstlisting}[style=c]
volatile char* vga = (volatile char*)0xB8000;
vga[0] = 'C';    /* character at row 0, column 0 */
vga[1] = 0x02;   /* attribute: background 0 (black), foreground 2 (green) */
\end{lstlisting}
\end{tryit}

\begin{warn}{Why this address works directly}
\hexa{B8000} lies below 4 MB, and the system maps the first 4 MB with \term{identity}: virtual
address \hexa{B8000} corresponds to physical \hexa{B8000}. That is why the driver can write there
without ceremony, even with paging enabled. If the mapping were different, translation would be
needed. You will see this in Part VIII.
\end{warn}

% ============================================================
\chapter{The text driver}

\begin{goals}
  \item See how a minimal driver is organised around a state structure
  \item Understand cursor handling and line wrapping
  \item Spot a latent divide-by-zero and understand why it is so dangerous
\end{goals}

\section{Terminal state}

A driver needs to remember things between calls: where the cursor is and what colour to write in.
All of that lives in a structure:

\filetag{lytlnybl/include/kernel/drivers/vga\_text.h}
\begin{lstlisting}[style=c]
typedef struct {
    size_t row;        /* current cursor row */
    size_t column;     /* current column */
    size_t width;      /* 80 */
    size_t height;     /* 25 */
    uint8_t color;     /* current attribute */
    uint16_t* buffer;  /* 0xB8000 */
} vga_text;
\end{lstlisting}

There is a single global instance, declared in \cod{kernel\_main.c} and used across the kernel with
\cod{extern}:

\begin{lstlisting}[style=c]
vga_text terminal;             /* in kernel_main.c */
extern vga_text terminal;      /* in every other file */
\end{lstlisting}

\begin{concept}{Passing state explicitly}
Notice that every function takes \cod{vga\_text* terminal} as its first argument, even though there
is only one instance. It is a deliberate design decision: it keeps the driver free of hidden state
and would allow several virtual terminals in future without touching a line of the logic. It is a
habit worth copying.
\end{concept}

\section{Initialising and clearing}

\filetag{lytlnybl/kernel/drivers/vga\_text.c}
\begin{lstlisting}[style=c]
void vga_text_init(vga_text* terminal) {
    terminal->row = 0;
    terminal->column = 0;
    vga_text_set_color(terminal, VGA_COLOR_WHITE, VGA_COLOR_RED);
    terminal->width = 80;
    terminal->height = 25;
    terminal->buffer = (uint16_t*)0xB8000;
    vga_text_clear(terminal);
}

void vga_text_clear(vga_text* terminal) {
    uint8_t color = terminal->color;
    uint16_t blank = ((uint16_t)color << 8) | ' ';
    for (size_t row = 0; row < terminal->height; row++) {
        for (size_t col = 0; col < terminal->width; col++) {
            size_t index = row * terminal->width + col;
            terminal->buffer[index] = blank;
        }
    }
    terminal->row = 0;
    terminal->column = 0;
}
\end{lstlisting}

``Clearing'' the screen means writing 2,000 spaces. There is no ``erase'' command: the hardware
always shows whatever is in those 4,000 bytes.

\section{Writing a string}

\filetag{lytlnybl/kernel/drivers/vga\_text.c}
\begin{lstlisting}[style=c]
void vga_text_write(vga_text* terminal, const char* string) {
    size_t pos = terminal->row * terminal->width + terminal->column;

    for (size_t i = 0; string[i] != '\0'; i++) {
        if (string[i] == '\n') {
            vga_text_writeline(terminal, "");
            continue;
        }
        vga_text_putchar(terminal, string[i]);
        pos++;
        terminal->row    = pos / terminal->width;
        terminal->column = pos % terminal->width;
        terminal->row    = terminal->row % terminal->height;
    }
}
\end{lstlisting}

The last three lines do all the positioning work with pure arithmetic:

\begin{itemize}
  \item \cod{pos / 80} gives the row: at position 80 it becomes 1, at 160 it becomes 2\ldots
  \item \cod{pos \% 80} gives the column within that row.
  \item \cod{row \% 25} makes row 24 wrap back to row 0.
\end{itemize}

\begin{warn}{Wrapping is not scrolling}
That last modulo makes the screen \emph{wrap}: on reaching the bottom it returns to the top and
writes over what was there. That is not what a real terminal does, which \emph{scrolls} everything
up by one line. The result is that long output ends up jumbled. Adding real scrolling is an
excellent exercise: you must move 24 rows with \cod{memcpy} and clear the last one.
\end{warn}

\section{A divide-by-zero waiting to happen}

Look at those three lines again. There are three divisions, and all three divide by a field of the
structure: \cod{terminal->width} and \cod{terminal->height}.

\begin{danger}{What happens if \cod{width} or \cod{height} is 0}
On x86, dividing by zero does not give ``infinity'' or a strange value: it \hi{raises
exception 0}, \emph{Divide By Zero}. The kernel handles it, and its handler does the most natural
thing in the world: print an error message.

But printing a message calls \cod{vga\_text\_writeline}, which calls \cod{vga\_text\_write}, which
divides by zero again. The exception fires again. And again. Each repetition consumes about 150
bytes of kernel stack, and \hi{never returns}.

After 107 rounds --- $107 \times 156 \approx 16$ KB, exactly the size of the kernel stack ---
the stack runs off the bottom of mapped memory, a page fault occurs, the
fault handler tries to use the same broken stack, a double fault occurs, and on the third failure
the processor gives up and resets the machine. What the user sees is a system that boots, gets to a
certain point, and reboots in a loop, with no message at all.

This is not hypothetical: Part XIII contains the real processor trace showing those 107 rounds, and
the step-by-step diagnosis down to the root cause. The cause was not in the video driver: it was
three subsystems away.
\end{danger}

\begin{concept}{The underlying lesson}
A kernel fault rarely shows up where its cause lives. Here, a rounding error in the physical memory
manager ended up causing a reboot loop in the video driver. Debugging an operating system means
thinking in \emph{chains} of causes, not about the point where it explodes. Part XIII teaches the
method.
\end{concept}

\section{Printing numbers}

There is no \cod{printf} here, so number-to-text conversion is done by hand:

\filetag{lytlnybl/kernel/drivers/vga\_text.c}
\begin{lstlisting}[style=c]
void vga_text_write_dec(vga_text* terminal, uint32_t value) {
    char buffer[10];
    size_t digits = 0;

    if (value == 0) {
        vga_text_putchar(terminal, '0');
        /* ... advance cursor ... */
        return;
    }

    while (value > 0) {
        buffer[digits++] = '0' + (value % 10);  /* least significant digit */
        value /= 10;
    }
    while (digits > 0) {                        /* print in reverse */
        vga_text_write_char(terminal, buffer[--digits]);
    }
}
\end{lstlisting}

The \cod{'0' + digit} trick works because ASCII digits are consecutive: \cod{'0'} is 48, \cod{'1'}
is 49, and so on. Adding the numeric value to the code of \cod{'0'} gives the right character.

And it must print in reverse because the algorithm extracts units first, then tens, and so on. Hence
storing them in an array and walking it backwards.

\begin{recipe}{everything the text driver can do}{ch18-vga-output}
A replacement \cod{kernel\_main.c} that exercises every function in the driver:
\cod{write} versus \cod{writeline}, decimal and hexadecimal, all sixteen colours, placing a
character at an exact cell with \cod{vga\_text\_put\_entry\_at}, and reading it back off the screen
with \cod{vga\_text\_getchar}.

Watch out for one real bug it exposes: \cod{vga\_text\_write\_dec(t, 0)} writes the \cod{'0'} with
\cod{putchar} but never advances the cursor, so the next thing printed covers it and the zero
vanishes. Every other value goes through a loop that does advance it.
\end{recipe}


\exercise{Write \cod{vga\_text\_write\_hex}, printing a \cod{uint32\_t} in hexadecimal with 8
digits and a \cod{0x} prefix. Hint: instead of \cod{\% 10} use \cod{\& 0xF} and shift 4 bits; to
turn a 0--15 value into a character you will need to distinguish whether it is below 10.}

\begin{summary}
With the video driver working, we can see what the system is doing --- which is the single most
important debugging tool in this environment. Next comes the most important topic in the course:
interrupts.
\end{summary}
